\chapter{联合熵、互信息与 KL 散度}

上一章只处理一个随机变量。本章研究两个变量之间的信息关系：知道 $Y$ 后，$X$ 还剩多少不确定性？两个变量共享多少统计信息？两个分布相差多少？

\section{联合分布、边缘分布与条件分布}

联合概率质量函数为 $p_{XY}(x,y)$。边缘分布由求和得到：
\[
p_X(x)=\sum_y p_{XY}(x,y),
\qquad
p_Y(y)=\sum_x p_{XY}(x,y).
\]
当 $p_Y(y)>0$ 时，条件分布为
\[
p_{X|Y}(x|y)=\frac{p_{XY}(x,y)}{p_Y(y)}.
\]
独立性等价于 $p_{XY}(x,y)=p_X(x)p_Y(y)$ 对所有 $(x,y)$ 成立。

\begin{definition}[联合熵与条件熵]
\begin{align*}
H(X,Y)&=-\sum_{x,y}p(x,y)\log_2p(x,y),\\
H(X|Y)&=\sum_y p(y)H(X|Y=y)\\
&=-\sum_{x,y}p(x,y)\log_2p(x|y).
\end{align*}
\end{definition}

\section{熵的链式法则}

\begin{theorem}[链式法则]
\[
H(X,Y)=H(Y)+H(X|Y)=H(X)+H(Y|X).
\]
更一般地，
\[
H(X_1,\dots,X_n)=\sum_{i=1}^n H(X_i|X_1,\dots,X_{i-1}).
\]
\end{theorem}

\begin{proof}
利用 $p(x,y)=p(y)p(x|y)$，有
\begin{align*}
H(X,Y)
&=-\sum_{x,y}p(x,y)\log_2[p(y)p(x|y)]\\
&=-\sum_{x,y}p(x,y)\log_2p(y)
  -\sum_{x,y}p(x,y)\log_2p(x|y)\\
&=H(Y)+H(X|Y).
\end{align*}
交换 $X,Y$ 得到另一种分解。多变量版本反复应用二变量链式法则即可。
\end{proof}

\begin{intuitionbox}
联合不确定性可以按顺序支付：先描述 $Y$，再在已知 $Y$ 的条件下描述 $X$。链式法则是编码、典型集和通信协议中最常用的“信息记账规则”。
\end{intuitionbox}

\section{KL 散度与 Gibbs 不等式}

\begin{definition}[KL 散度]
对有限集合上的分布 $P,Q$，若 $P(x)>0$ 时均有 $Q(x)>0$，定义
\[
D_{\mathrm{KL}}(P\|Q)
=\sum_x P(x)\log_2\frac{P(x)}{Q(x)}.
\]
若存在 $P(x)>0$ 但 $Q(x)=0$，则散度为 $+\infty$。
\end{definition}

\begin{theorem}[Gibbs 不等式]
\[
D_{\mathrm{KL}}(P\|Q)\ge 0,
\]
等号当且仅当 $P=Q$（在 $P$ 的支持上）。
\end{theorem}

\begin{proof}
使用自然对数不等式 $\ln t\le t-1$：
\begin{align*}
D_{\mathrm{KL}}(P\|Q)
&=-\frac1{\ln2}\sum_xP(x)\ln\frac{Q(x)}{P(x)}\\
&\ge -\frac1{\ln2}\sum_xP(x)
\left(\frac{Q(x)}{P(x)}-1\right)\\
&=-\frac1{\ln2}\left(\sum_xQ(x)-\sum_xP(x)\right)=0.
\end{align*}
等号要求 $Q(x)/P(x)=1$ 对所有 $P(x)>0$ 成立。
\end{proof}

\subsection{交叉熵}

交叉熵定义为
\[
H(P,Q)=-\sum_xP(x)\log_2Q(x).
\]
它满足
\[
H(P,Q)=H(P)+D_{\mathrm{KL}}(P\|Q).
\]
当真实分布固定时，最小化交叉熵等价于最小化 $D_{\mathrm{KL}}(P\|Q)$。这解释了分类模型中的交叉熵目标，但不意味着交叉熵下降自动带来推理、证明或算法能力。

\begin{warningbox}
KL 散度不是距离：通常 $D(P\|Q)\ne D(Q\|P)$，也不满足三角不等式。书写时不能省略方向。
\end{warningbox}

\begin{example}[KL 的不对称性]
取 $P=(1/2,1/2)$、$Q=(3/4,1/4)$，则
\begin{align*}
D(P\|Q)&=\tfrac12\log_2\tfrac23+\tfrac12\log_2 2\approx0.2075,\\
D(Q\|P)&=\tfrac34\log_2\tfrac32+\tfrac14\log_2\tfrac12\approx0.1887.
\end{align*}
两者不同。
\end{example}

\section{互信息}

\begin{definition}[互信息]
\[
I(X;Y)=D_{\mathrm{KL}}(P_{XY}\|P_XP_Y)
=\sum_{x,y}p(x,y)\log_2\frac{p(x,y)}{p(x)p(y)}.
\]
\end{definition}

由 KL 非负性，$I(X;Y)\ge0$；等号当且仅当 $X,Y$ 独立。

把联合分布分解为 $p(x,y)=p(y)p(x|y)$，可得四个等价表达：
\begin{align*}
I(X;Y)
&=H(X)-H(X|Y)\\
&=H(Y)-H(Y|X)\\
&=H(X)+H(Y)-H(X,Y)\\
&=D_{\mathrm{KL}}(P_{XY}\|P_XP_Y).
\end{align*}

因此，互信息既可以解释为“观察 $Y$ 后 $X$ 的不确定性平均减少多少”，也可以解释为“联合分布偏离独立分布多远”。

\begin{example}[带噪复制]
设 $X\sim\mathrm{Bernoulli}(1/2)$，$N\sim\mathrm{Bernoulli}(q)$ 且独立，令 $Y=X\oplus N$。则 $Y$ 仍均匀，
\[
H(Y)=1,\qquad H(Y|X)=H(N)=h_2(q),
\]
所以
\[
I(X;Y)=1-h_2(q).
\]
当 $q=0$ 时完全保留 1 bit；当 $q=1/2$ 时输出与输入独立，互信息为 0。
\end{example}

\section{条件互信息与链式法则}

\begin{definition}[条件互信息]
\[
I(X;Y|Z)=H(X|Z)-H(X|Y,Z).
\]
\end{definition}
它是对每个 $Z=z$ 的互信息取平均，因此非负。互信息链式法则为
\[
I(X;Y,Z)=I(X;Z)+I(X;Y|Z).
\]
进一步，
\[
I(X;Y_1,\dots,Y_n)
=\sum_{i=1}^n I(X;Y_i|Y_1,\dots,Y_{i-1}).
\]

\section{条件化减少熵与数据处理不等式}

由 $I(X;Y)=H(X)-H(X|Y)\ge0$，得到
\[
H(X|Y)\le H(X).
\]
这是一条平均意义的结论；对某个特定 $Y=y$，条件熵可以高于原熵。

\begin{theorem}[数据处理不等式]
若 $X\to Y\to Z$ 构成 Markov 链，即给定 $Y$ 后 $Z$ 与 $X$ 条件独立，则
\[
I(X;Z)\le I(X;Y).
\]
\end{theorem}

\begin{proof}[证明骨架]
一方面，由链式法则
\[
I(X;Y,Z)=I(X;Y)+I(X;Z|Y)=I(X;Y),
\]
其中 Markov 条件给出 $I(X;Z|Y)=0$。另一方面，
\[
I(X;Y,Z)=I(X;Z)+I(X;Y|Z)\ge I(X;Z).
\]
合并即得结论。
\end{proof}

\begin{intuitionbox}
对数据做后处理不能凭空增加它与原始变量之间的统计信息。后处理可以重新组织、压缩或突出某些信息，但若只观察 $Z=f(Y)$，关于 $X$ 的总互信息不会超过观察完整 $Y$。
\end{intuitionbox}

\section{一个联合分布例子}

令联合分布为
\[
\begin{array}{c|cc}
 &Y=0&Y=1\\\hline
X=0&3/8&1/8\\
X=1&1/8&3/8
\end{array}
\]
两个边缘分布都均匀，所以 $H(X)=H(Y)=1$。给定 $Y$ 后，$X=Y$ 的概率为 $3/4$，故
\[
H(X|Y)=h_2(1/4)\approx0.811,
\qquad
I(X;Y)\approx0.189\bits.
\]
这个数值等于一个交叉概率 $q=1/4$ 的二元对称信道所保留的信息。

\section{小结与验收}

\begin{checkpointbox}
\begin{enumerate}
  \item 从 $p(x,y)=p(y)p(x|y)$ 推导熵的链式法则；
  \item 说明 KL 散度非负的关键不等式和等号条件；
  \item 写出互信息的四种等价表达，并解释每一种视角；
  \item 对上面的 $2\times2$ 联合分布手算 $H(X|Y)$ 和 $I(X;Y)$；
  \item 用一句话说明数据处理不等式依赖的 Markov 假设。
\end{enumerate}
\end{checkpointbox}
